\section{练习}

\begin{exercise}
在 \(\RR^3\) 上取
\[
e_1=\partial_x+y\partial_z,
\qquad
e_2=\partial_y,
\qquad
e_3=\partial_z.
\]
求对偶余标架与全部非零结构函数，并直接验证\eqnrefs{eq:math-dg-frame-structure-equation}。
\end{exercise}

\begin{exercise}
设 \(e'_a=e_b\Lambda^b{}_a\)。从 \(\nabla e'_a\) 的定义重新推导联络换标架律\eqnrefs{eq:math-dg-connection-gauge-transform}，并据此证明曲率满足\eqnrefs{eq:math-dg-curvature-gauge-transform}。
\end{exercise}

\begin{exercise}
从 Cartan 第一、第二结构方程出发，不使用分量公式，完整推导两条 Bianchi 恒等\eqnrefs{eq:math-dg-bianchi-first}与\eqnrefs{eq:math-dg-bianchi-second}。
\end{exercise}

\begin{exercise}
对二维度量
\[
g=\dd r^2+f(r)^2\dd\phi^2,
\qquad f(r)>0,
\]
取正交余标架 \(\theta^1=\dd r\)、\(\theta^2=f(r)\dd\phi\)。先由无挠条件求 Levi--Civita 联络一形式，再由第二结构方程证明 Gaussian 曲率 为
\[
K=-\frac{f''(r)}{f(r)}.
\]
分别代入 \(f(r)=r\)、\(\sin r\)、\(\sinh r\) 检查平面、单位球面与双曲平面的曲率。
\end{exercise}

\begin{exercise}
设 \(X\) 为 Killing 场。利用\eqnrefs{eq:math-dg-synthesis-hodge-commute}与 \([\mathcal L_X,\dd]=0\) 证明 \([\mathcal L_X,\Delta]=0\)。说明为什么这个结论意味着每个 Hodge--拉普拉斯算子 特征空间在 Killing 流下保持。
\end{exercise}

\begin{exercise}
设 \(M\) 为紧致定向无边界 Riemann 流形。若闭形式 \(\alpha\) 的 Hodge 分解为
\[
\alpha=\dd\beta+\delta\gamma+\alpha_{\mathcal H},
\]
不用直接引用正文结论，证明 \(\delta\gamma=0\)，从而说明 \([\alpha]=[\alpha_{\mathcal H}]\)。再说明为什么 调和 分量 \(\alpha_{\mathcal H}\) 唯一而 \(\beta\) 不必唯一。
\end{exercise}

\begin{exercise}
在 \(\RR^3\) 中考虑一形式
\[
\alpha=\dd z-x\,\dd y.
\]
求 \(\alpha\wedge\dd\alpha\)，据此判断 \(\ker\alpha\) 是否可积。再自行构造一个非零一形式 \(\widetilde\alpha\) 使 \(\widetilde\alpha\wedge\dd\widetilde\alpha=0\)，并写出其局部积分曲面。
\end{exercise}

\begin{exercise}
设 \(g(t)\) 为度量族，记 \(\dot g=\partial_tg\)。从 \(g^{ik}g_{kj}=\delta^i{}_j\) 与 Jacobi 行列式恒等式分别推导\eqnrefs{eq:math-dg-inverse-metric-variation}和\eqnrefs{eq:math-dg-volume-variation}。然后令 \(\dot g=-2\operatorname{Ric}\)，推出\eqnrefs{eq:math-dg-ricci-flow-volume}。
\end{exercise}

\begin{exercise}
设 \(\alpha_t\) 满足纯输运方程
\[
\partial_t\alpha_t+\mathcal L_{V_t}\alpha_t=0.
\]
由\eqnrefs{eq:math-dg-transport-identity}证明 \(\Phi_{t,s}^*\alpha_t=\alpha_s\)。若 \(\alpha_s\) 闭，证明所有时刻都闭；若 \(\alpha_s\) 恰当，证明所有时刻都恰当。
\end{exercise}

\begin{exercise}
从 Gauss--Weingarten 分解\eqnrefs{eq:math-dg-gauss-formula-general}与\eqnrefs{eq:math-dg-weingarten-formula-general}重新计算 \(\overline R(X,Y)Z\) 的切向和法向部分，推导一般余维 Gauss 与 Codazzi 方程。然后计算 \(\overline R(X,Y)\xi\) 的法向部分，核对 Ricci 方程中交换子 \([A_\xi,A_\eta]\) 的次序与符号。
\end{exercise}

\begin{exercise}
对单位球面取外法向，按照本书约定 \(\overline\nabla_Xn=-AX\) 求 形状算符、两个主曲率、Gaussian 曲率 与有向平均曲率 \(\mathsf H\)。反转法向后重新计算，并说明哪些量改变符号、哪些量保持不变。
\end{exercise}
